\section*{September 6, 2026: Recovering the construction data's derivation}
\textit{Codex draft. The restoration remains in progress on the sales\_sample\_exploration branch.}

The construction decisions and much of their supporting code survive in the
research archive. The earlier claim that the central cleaning could not be
traced was too strong. Production had become disconnected from that derivation:
it started with a saved analysis dataset rather than rebuilding the projects
from source records and recorded decisions. Jacob authorized restoring that
connection on this branch. Repeating human reviews each time the paper compiles
is unnecessary; their recorded decisions must enter the build as inputs.

Replaying the recovered code has reproduced substantial parts of the earlier
cleaning. The residential review stage reproduces all 184 source dispositions
and its 139 retained projects. The commercial decision stage reproduces all
1,128 source dispositions and 815 retained projects. Historical parcel matching
reproduces the saved coverage and predecessor decisions, including the geometry
and attributes of 11,794 source parcels and 3,770 selected predecessor
geometries. These checks use preserved upstream inputs where their producers
have not yet been restored. An earlier partial chain also reproduced the
8,648-row final construction export. None of these comparisons establishes a
complete rebuild from raw data with unchanged estimates.

Two confirmed programming errors affected the preservation of conflicting
source information. The historical coordinate screen counted distinct unit and
building-area values after aggregation had already replaced them. Counting the
original values first reduces the accepted coordinate file from 391 to 387
rows. The four removed locations are outside the 500-foot sample. Separately,
the commercial summary discarded lists of competing source values after
collapsing them to a single value. Preserving those lists restores evidence
about conflicting years, units, and areas; the commercial candidate decisions
remain unchanged in the replay. These corrections are in the restored code.
The production construction export has not been replaced.

The exact original full residential Assessor download remains unlocated.
Jacob does not know of another backup. A fresh, pinned download contains every
previously selected residential PIN and 199 additional PINs. Its effects require
more than a row-count comparison. With the preserved candidate inventory, it
makes two previously unresolved parcel groups eligible for mechanical retention
because they now have a complete 2026 assessment snapshot. Both projects already
enter the frozen final dataset through later resolution. One lies 419 feet from
its ward boundary: the later snapshot reports 2,976 square feet of land and
3,115 square feet of building area, whereas the existing resolution uses 1,440
and 3,114, respectively. The earlier resolution comes from selected 2022
building-card evidence. This difference has not been adopted. The existing
review-coverage check rejects a changed review cohort; it must not be removed
merely to make the new source pass.

Zoning replay also revealed that the grouping function called an absent map
match ``Other.'' Preserving missingness instead changes one project's 2012
snapshot evidence. Its final Downtown classification still comes from validated
component evidence. In the replay, no final zoning group or assignment source
changes. The missing-value correction is in both restored zoning stages.

The remaining work is to restore the outstanding parcel, building, permit,
and historical zoning producers, reconcile the replacement source with the
recorded resolutions, and run the resulting chain through the paper. Source
extracts need explicit query coverage so that an unqueried location cannot
be classified as having no match. Jacob also authorized moving the shared
files required by this pipeline to \path{tasks/shared/code/} and updating their
callers. A separate Make 3.81 experiment verified recovery when either member
of a two-output producer is deleted, including parallel execution. The full
paper build still needs verification; changing shared include paths must not
silently refresh mutable source datasets. The parcel-universe task now replays
its preserved source snapshot by default, with a separate Make entry point for
a deliberate API refresh. Its 882,741-row consumer dataset reproduces the
existing file byte for byte. Both the source and rewritten data have unique,
nonmissing parcel identifiers. Other mutable downloads still need this vintage
check before running the full paper build.

A separate comparison raises a land-denominator question for retained projects
with several building cards. Among 248 such projects in the frozen paper data,
89 have a recorded successor match for every target card. Four projects within
500 feet have complete matches that reproduce their building totals, yet their
parent land area exceeds the matched successor lots' combined area by more than
a factor of two. This is a diagnostic comparison, not a new exclusion rule.

The largest difference concerns the seven-card Park Place Homes parent,
PIN 19111200130000. The paper data assign its seven units and 13,753 square feet
of building area to 307,739 square feet of land. Its seven matched successor
lots total 21,069 square feet, a factor of 14.6 smaller. All seven successor rows
are absent from the final dataset, so this comparison does not demonstrate
double counting. The recorded review supports the card inventory; a later web
review identifies the development but explicitly does not independently verify
the parent's unit count. Subdivision and common land may explain some of the
land difference. The outstanding question is whether the numerator and
denominator describe the same construction episode. No replacement land area
or sample exclusion has been adopted.

The follow-up review checked all four parent land values against the selected
historical Assessor cards and mapped parcels. The 13-home Calumet parent's
27,738 square feet agree with its 2008 polygon; all 13 matched homes lie inside,
with no additional reported residential units. Keeping that denominator is
supported because the smaller individual lots exclude additional site land.
The other three cases remain open. Fletcher's parent covers five homes while
its numerator counts two, and the other three homes remain separate observations.
Park Place includes a much larger tract and later development phases. The
29-home Eastgate parent's area is confirmed by its 2008 polygon, but the site
contains six additional homes in the recorded successor evidence. Allocating
land or changing episode definitions is necessary before clearing those density
values. These are review recommendations; no production values were changed.
The case evidence and maps are generated by the release audit's
\path{review_denominator_cases.R}; the detailed interpretation is in
\path{denominator_review.md} in that audit directory.

\small
The permit follow-up makes the reported years less reliable for these sites.
Fletcher's selected homes and the two Campbell Avenue homes received new-construction
permits on December 1, 2006. The permits identify components of larger building
groups, so the five homes inside the tax parent are not necessarily one complete
building. At Eastgate, the five southern homes received a new-construction permit
on August 10, 2006, although four successor records report year built 2005.
Permit issuance is not completion evidence, but these records do not support
classifying those four homes as established pre-study housing. Park Place's
seven matched successors span the subareas in the 2013 development drawing;
the drawing's separate seven-unit subarea is therefore not a valid substitute.
No new denominator or year has been adopted. The question of excluding unresolved
density observations versus imputing shared land remains with Jacob.

A separate restoration reproduced the recorded evidence for construction episodes
linked across parcel identifiers: 123 project records and 596 assessment-year
episode records match the preserved files byte for byte. The restored Make rules
regenerated a deleted secondary output during a parallel build, then ran no
producer on the unchanged next build. This check uses the preserved review bundle
and assessment history, so it does not establish reconciliation of the replacement
source vintage or completion of the full upstream chain.

\textbf{Reproduction note.}
The recovered code is from research-archive commit
\nolinkurl{010a1f8497c1}.
The new production task is \path{tasks/new_construction_cleaning/}; its default
build is incomplete while upstream producers remain missing. Recorded stage
comparisons and the four affected tieback candidates are documented in
the release audit files \path{restored_stage_checks.csv} and
\path{current_assessor_temporal_candidate_changes.csv}.
Source vintages and hashes are in \path{data_raw/construction_review/}.
The land comparison is generated by the release audit's
\path{code/review_multicard_land_scope.R} from the frozen paper data and the
recorded review baseline in that audit's \path{reference/} directory. These
baseline files are computed historical review products, not raw data or a
substitute for restoring their production derivation.
The experiments were run locally; they do not establish execution on a cluster.
Compile this entry with \texttt{make} in \path{logbook/}.
\normalsize
